% ------------------------------------------------------------------ 4. Despliegue
\seccionnueva
\section{Estructura de despliegue}
\label{sec:despliegue}

La vista de despliegue asigna los componentes de la sección~\ref{sec:cc} a nodos de ejecución.
Hay dos tipos de nodo: el \textbf{host cliente} (el computador de la persona usuaria, con el
navegador y el ejecutable de escritorio) y el \textbf{host Docker}, donde Docker Compose ejecuta
cada servicio, el broker y cada base de datos en su propio contenedor, conectados por una red
interna. El único puerto publicado hacia el host cliente es el del gateway.

\begin{nota}[Frontera explícita]
El cliente local (Electron) \textbf{no corre en Docker}: se empaqueta como ejecutable y se
instala en el host cliente. Un contenedor puede usarse para compilarlo, pero la interfaz gráfica
se ejecuta fuera de Docker y habla con el gateway igual que el navegador.
\end{nota}

\begin{landscape}
\begin{figure}[p]
\centering
\begin{tikzpicture}[x=1cm,y=1cm]
  \tikzset{
    ctr/.style={container,minimum width=2.75cm,text width=2.55cm},
    ctrdb/.style={container,draw=cbGreen!70!black,fill=cbSoftGreen,minimum width=2.75cm,text width=2.55cm},
    ctrmq/.style={container,draw=cbAmber,fill=cbSoftAmber,minimum width=2.75cm,text width=2.55cm},
    ctrweb/.style={container,draw=cbViolet!80!black,fill=cbSoftViolet,minimum width=2.75cm,text width=2.55cm},
    vol/.style={volume},
  }

  % ---------------- Host cliente
  \node[ctrweb,minimum height=16mm] (browser) at (-10.3,-2.2)
    {\textbf{Navegador web}\\HTML del SSR\\+ JS hidratado};
  \node[ctrweb,minimum height=10mm] (main) at (-10.3,-6.0) {\textbf{Proceso principal}\\Node · token · HTTP/WS};
  \node[ctrweb,minimum height=10mm] (renderer) at (-10.3,-8.1) {\textbf{Renderer}\\React · \texttt{ui/} · IPC};
  \begin{scope}[on background layer]
    % Primero el nodo exterior y luego el interior, para que el interior quede visible.
    \node[node3d,fit=(browser)(main)(renderer),inner sep=22pt,
      label={[font=\sffamily\small\bfseries,anchor=north west]north west:«device» Host cliente}] (client) {};
    \node[draw=cbViolet!70,fill=cbSoftViolet!35,rounded corners=3pt,fit=(main)(renderer),inner sep=8pt,
      label={[font=\sffamily\scriptsize\bfseries,text=cbViolet]above:CollaBoard Desktop (Electron)}] (electron) {};
  \end{scope}
  \node[font=\sffamily\scriptsize,text=cbMuted,anchor=north] at (client.south) {Windows 11 · fuera de Docker};

  % ---------------- Host Docker
  % Ingreso
  \node[ctrmq,minimum height=13mm] (gw) at (-4.3,-1.6) {\textbf{gateway}\\nginx:1.27\\{\tiny :80 · publicado}};
  \node[ctrweb,minimum height=13mm] (web) at (-0.9,-1.6) {\textbf{web-ssr}\\node:22 · Next.js\\{\tiny :3000 · /healthz}};
  % Lógica
  \node[ctr] (auth)   at (-4.3,-4.4) {\textbf{auth-api}\\Go\\{\tiny :8001}};
  \node[ctr] (collab) at (-0.9,-4.4) {\textbf{colaboracion-api}\\Python\\{\tiny :8000}};
  \node[ctr] (canvas) at (2.5,-4.4)  {\textbf{pizarra-api}\\Python\\{\tiny :8000 · WS}};
  \node[ctr] (chat)   at (5.9,-4.4)  {\textbf{chat-api}\\Java\\{\tiny HTTP + WS}};
  \node[ctr] (pixels) at (2.5,-1.6)  {\textbf{pixels-api}\\Python\\{\tiny :8000 · WS}};
  \node[ctr] (notif)  at (5.9,-1.6)  {\textbf{notification-api}\\C\#\\{\tiny HTTP}};
  \node[ctrmq] (mq)   at (9.3,-1.6)  {\textbf{rabbitmq}\\rabbitmq:3\\{\tiny :5672 · DLQ}};
  \node[ctr] (worker) at (9.3,-4.4)  {\textbf{message-worker}\\C\#\\{\tiny consumidor}};
  % Datos
  \node[ctrdb] (userdb)  at (-4.3,-7.3) {\textbf{user-db}\\postgres:16};
  \node[ctrdb] (boarddb) at (-0.9,-7.3) {\textbf{board-db}\\postgres:16};
  \node[ctrdb] (canvasdb) at (2.5,-7.3) {\textbf{canvas-db}\\postgres:16};
  \node[ctrdb] (mongo)   at (5.9,-7.3)  {\textbf{pixels-nosql-db}\\mongo:7};
  \node[ctrdb] (notifdb) at (9.3,-7.3)  {\textbf{notification-db}\\postgres:16};
  % Volúmenes
  \node[vol] (v1) at (-4.3,-9.35) {user\_data};
  \node[vol] (v2) at (-0.9,-9.35) {board\_data};
  \node[vol] (v3) at (2.5,-9.35)  {canvas\_data};
  \node[vol] (v4) at (5.9,-9.35)  {mongo\_data};
  \node[vol] (v5) at (9.3,-9.35)  {notif\_data};
  \foreach \d/\v in {userdb/v1,boarddb/v2,canvasdb/v3,mongo/v4,notifdb/v5}
    \draw[thick,cbGreen!70!black,densely dotted] (\d) -- (\v);

  \begin{scope}[on background layer]
    \node[node3d,fit=(gw)(mq)(v1)(v5),inner sep=22pt,
      label={[font=\sffamily\small\bfseries,anchor=north west]north west:«device» Host Docker (Docker Engine + Compose)}] (docker) {};
    \node[draw=cbIndigo!50,dashed,rounded corners=4pt,fill=cbGray,fit=(gw)(mq)(worker)(userdb)(notifdb),inner sep=10pt,
      label={[font=\sffamily\scriptsize\bfseries,text=cbIndigoDark,anchor=north east]south east:red interna \texttt{collaboard-network} (bridge)}] (net) {};
  \end{scope}

  % ---------------- Comunicación entre nodos
  \draw[http] (browser.east) -- node[lbl,above,pos=0.45] {HTTP + WS} (gw.west);
  \draw[http] (main.east) to[out=0,in=225] node[lbl,below,pos=0.35] {HTTP + WS (Bearer)} (gw.south west);
  \draw[<->,thick,cbViolet] (main) -- node[lbl,right] {IPC} (renderer);
  \draw[http] (gw) -- (web);
  \node[font=\sffamily\scriptsize,text=cbMuted,align=left,anchor=north west] at ([yshift=-4pt]docker.south west)
    {Sin puertos publicados para APIs, broker ni bases: solo son accesibles dentro de la red interna.\\
     Las líneas punteadas verdes son volúmenes Docker con nombre (los datos sobreviven a reinicios).};
\end{tikzpicture}
\caption{Vista de despliegue: host cliente (navegador y Electron) y host Docker con quince
contenedores en una red interna; solo el gateway publica un puerto.}
\label{fig:despliegue}
\end{figure}
\end{landscape}

\subsection{Elementos de despliegue}

\begin{longtable}{@{}P{3.2cm} P{2.9cm} P{2.3cm} P{2.4cm} P{3.6cm}@{}}
  \encabezado \hd{Contenedor} & \hd{Imagen / runtime} & \hd{Puerto} & \hd{Volumen} & \hd{Healthcheck} \\
  \endfirsthead
  \encabezado \hd{Contenedor} & \hd{Imagen / runtime} & \hd{Puerto} & \hd{Volumen} & \hd{Healthcheck} \\
  \endhead
  gateway & nginx:1.27-alpine & publicado & — & \code{GET /healthz} (a web-ssr) \\
  \rowcolor{cbGray} web-ssr & node:22-alpine (standalone) & 3000 interno & — & \code{GET /healthz} \\
  auth-api & Go 1.23 / alpine & 8001 interno & — & \code{GET /health} \\
  \rowcolor{cbGray} colaboracion-api & python:3.12-slim & 8000 interno & — & \code{GET /health} \\
  pizarra-api & python:3.12-slim & 8000 interno & — & \code{GET /health} \\
  \rowcolor{cbGray} chat-api & JVM & interno & — & por definir \\
  pixels-api & python:3.12-slim & 8000 interno & — & \code{GET /api/health} \\
  \rowcolor{cbGray} notification-api & .NET & interno & — & por definir \\
  message-worker & .NET & — & — & proceso vivo \\
  \rowcolor{cbGray} rabbitmq & rabbitmq:3-management & 5672 interno & rabbit\_data & \code{diagnostics ping} \\
  user-db, board-db, canvas-db, notification-db & postgres:16-alpine & 5432 interno & uno por base & \code{pg\_isready} \\
  \rowcolor{cbGray} pixels-nosql-db & mongo:7 & 27017 interno & mongo\_data & \code{mongosh ping} \\
  \bottomrule
  \caption{Contenedores del host Docker.}
  \label{tab:contenedores}
\end{longtable}

\subsection{Patrones de despliegue}

\begin{itemize}
  \item \textbf{Un contenedor por componente}: cada servicio, el \textit{worker}, el broker y cada base se construyen, escalan y reinician por separado.
  \item \textbf{Punto de entrada único}: solo el gateway publica un puerto; la web y las APIs comparten origen, lo que habilita la cookie de sesión y evita CORS.
  \item \textbf{Red interna aislada}: los servicios se encuentran por su nombre DNS en Compose (p.\,ej. \code{colaboracion-api}); las bases no son alcanzables desde el host cliente.
  \item \textbf{Datos en volúmenes con nombre}: reiniciar o reconstruir un contenedor no borra datos; las pruebas de reinicio se hacen sin \code{docker compose down -v}.
  \item \textbf{Healthchecks y \code{depends\_on: service\_healthy}}: los servicios arrancan cuando sus dependencias responden.
  \item \textbf{Configuración por entorno}: la misma imagen sirve en cualquier despliegue (la URL interna del gateway que usa la web se lee en tiempo de ejecución).
  \item \textbf{Cliente pesado empaquetado}: el ejecutable de Electron se distribuye aparte y se configura con la URL del gateway.
\end{itemize}

\begin{pendiente}{Sebastian (vista de despliegue)}
Ajustar nombres de servicios, puertos, volúmenes e imágenes al \code{docker-compose.yml} final y
confirmar que la figura y la tabla coinciden con él. Hoy existe un Compose de desarrollo para la
web (\code{web-ssr/compose.dev.yml}) que publica el gateway de pruebas en el puerto 8088.
\end{pendiente}
